% from ega4-2-fr.tex
\oldpage[IV-2]{111}

Mais comme le foncteur $f^*$ commute aux limites inductives
$(0_{\mathrm I},4.3.2)$, cela donne par définition l'isomorphisme
(5.9.1.1) cherché.

\begin{corollary}[5.9.5]
\label{IV.2.s5.9.5-fr}
Sous les hypothèses de (5.9.4), si
$\mathcal H^0_{X/Z}(\mathcal F)$ est cohérent, il en est de même de
$\mathcal H^0_{X'/Z'}(f^*(\mathcal F))$. La réciproque est vraie lorsque
$f$ est un morphisme fidèlement plat et quasi-compact.
\end{corollary}

La première assertion résulte de (5.9.4.1) et de $(0_{\mathrm I},5.3.11)$;
la seconde équivaut à dire que si
$\mathcal H^0_{X'/Z'}(f^*(\mathcal F))$ est de type fini, il en est de même
de $\mathcal H^0_{X/Z}(\mathcal F)$; cela résulte de (5.9.4.1) et de
(2.5.2).

\begin{corollary}[5.9.6]
\label{IV.2.s5.9.6-fr}
Soient $X$ un préschéma localement noethérien, $Z$ une partie de $X$ stable
par spécialisation, $\mathcal F$ un $\mathcal O_X$-Module quasi-cohérent.
Pour tout $x\in X$, posons $X_x=\operatorname{Spec}(\mathcal O_x)$,
$Z_x=Z\cap X_x$; on a un isomorphisme canonique fonctoriel
\begin{equation}
(\mathcal H^0_{X/Z}(\mathcal F))_x
\xrightarrow{\sim}\mathcal H^0_{X_x/Z_x}(\widetilde{\mathcal F}_x).
\tag{5.9.6.1}\label{IV.2.s5.9.6.1-fr}
\end{equation}
\end{corollary}

Il suffit d'appliquer (5.9.4) au morphisme canonique $X_x\to X$, qui est
plat, et de tenir compte de (I, 1.6.5).

\begin{env}[5.9.7]
\label{IV.2.s5.9.7-fr}
Avec les notations de (5.9.1), on a pour tout $\alpha$ un homomorphisme
canonique fonctoriel
\[
\mathcal F\longrightarrow(i_\alpha)_*(\mathcal F|U_\alpha)
\quad(0_{\mathrm I},4.4.3.2),
\]
et ces homomorphismes forment un système inductif; par passage à la limite
inductive, on en déduit donc un homomorphisme canonique fonctoriel
\begin{equation}
\rho_{X/Z}:\mathcal F\longrightarrow\mathcal H^0_{X/Z}(\mathcal F).
\tag{5.9.7.1}\label{IV.2.s5.9.7.1-fr}
\end{equation}
\end{env}

\begin{proposition}[5.9.8]
\EGAReaderTarget{ega4.c04.prop.5-9-8}{IV.2.s5.9.8-fr}
Soient $X$ un préschéma localement noethérien, $Z$ une partie de $X$ stable
par spécialisation, $\mathcal F$ un $\mathcal O_X$-Module. Les propriétés
suivantes sont équivalentes:
\begin{enumerate}
\item[a)] L'homomorphisme $\rho_{X/Z}$ (5.9.7.1) est injectif
(resp. bijectif).
\item[b)] Pour tout ouvert noethérien $V$ de $X$, l'homomorphisme
\[
(\rho_{X/Z})_V:\Gamma(V,\mathcal F)\longrightarrow
\varinjlim_\alpha\Gamma(V\cap U_\alpha,\mathcal F)
\]
est injectif (resp. bijectif).
\item[a')] Pour toute partie fermée $T\subset Z$ de $X$, l'homomorphisme
canonique $(0_{\mathrm I},4.4.3.2)$
\[
\mathcal F\longrightarrow i_*(\mathcal F|X-T)
\]
(où $i:X-T\to X$ est l'injection canonique) est injectif
(resp. bijectif).
\item[b')] Pour toute partie fermée $T\subset Z$ de $X$ et tout ouvert
noethérien $V$ de $X$, l'homomorphisme de restriction
\[
\Gamma(V,\mathcal F)\longrightarrow
\Gamma(V\cap(X-T),\mathcal F)
\]
est injectif (resp. bijectif).
\end{enumerate}
\end{proposition}

Compte tenu de (5.9.1.2), l'équivalence de a) et b) (resp. a') et b'))
résulte de
